% Theory / Theoretical Analysis Section Draft
% For: Recursive Generation Depth Estimation Paper

\section{Theoretical Analysis}
\label{sec:theory}

This section establishes the theoretical foundations of recursive generation depth estimation. We first connect our feature design to the distributional dynamics of iterative model retraining, then derive an upper bound on the Effective Detection Depth (EDD), and finally validate both against experimental observations.

\subsection{Distributional Dynamics of Recursive Generation}
\label{sec:theory:dynamics}

We model recursive generation as an iterative mapping on distributions. Let $\mathcal{D}_0$ denote the original (human-written) data distribution. At each generation step $k \geq 1$, a language model $M_k$ is trained on $\mathcal{D}_{k-1}$, and $\mathcal{D}_k$ is produced by sampling from $M_k$:
\begin{align}
    M_k &= \textsc{Train}(\mathcal{D}_{k-1}), \label{eq:train} \\
    \mathcal{D}_k &= \textsc{Sample}(M_k). \label{eq:sample}
\end{align}
This defines a Markov chain on the space of distributions, $\mathcal{D}_0 \to \mathcal{D}_1 \to \mathcal{D}_2 \to \cdots$, whose convergence properties determine both the nature of model collapse and the feasibility of depth estimation.

\begin{theorem}[Tail Contraction under Recursive Generation]
\label{thm:tail}
Let $p_k$ denote the probability density of $\mathcal{D}_k$. Under recursive generation with finite samples at each step, the tail mass
\begin{equation}
    T_k(\alpha) \triangleq \int_{\{x : p_0(x) < \alpha\}} p_k(x) \, dx
\end{equation}
satisfies $T_{k+1}(\alpha) \leq T_k(\alpha)$ for all $\alpha > 0$ and $k \geq 0$, with convergence $T_k(\alpha) \to 0$ as $k \to \infty$.
\end{theorem}

This result follows from the analysis of \citet{shumailov2024collapse}, who proved that for Gaussian models under recursive estimation, the variance $\sigma_k^2$ converges to zero, so the distribution degenerates toward a point mass. The mechanism is straightforward: finite-sample estimation at each step fails to capture low-probability events in the tails of $\mathcal{D}_{k-1}$, and this information loss compounds across iterations. \citet{dohmatob2024demystified} further showed that in high-dimensional regression settings, the test risk after $k$ iterations of retraining grows as
\begin{equation}
    R_k = R_0 + k \cdot \Delta_{\text{collapse}},
    \label{eq:linear_growth}
\end{equation}
where $\Delta_{\text{collapse}} > 0$ depends on the data covariance spectrum, regularization strength, and label noise level. In the overparameterized regime, this degradation accelerates to exponential rates.

\subsection{Why Perplexity Features Distinguish Depth}
\label{sec:theory:features}

We now explain why the two core features in our detection pipeline, the 90th-percentile perplexity (\texttt{p90\_ppl}) and the variance of token-level surprisal (\texttt{var\_surprisal}), provide monotonically decreasing signals as generation depth $k$ increases.

\paragraph{Tail perplexity (\texttt{p90\_ppl}).}
Perplexity at the 90th percentile captures the behavior of a reference model on the most surprising tokens in a text. These tokens correspond to rare constructions, unusual word choices, and low-frequency collocations that reside in the tail of the original distribution $\mathcal{D}_0$. By Theorem~\ref{thm:tail}, the tail mass $T_k(\alpha)$ decreases monotonically with depth $k$. As $\mathcal{D}_k$ loses these rare elements, texts sampled from deeper generations contain fewer high-surprisal tokens, and \texttt{p90\_ppl} decreases accordingly.

More precisely, let $s(x_t \mid x_{<t})$ denote the token-level surprisal assigned by a fixed reference model. For text $\mathbf{x} \sim \mathcal{D}_k$, define the empirical surprisal distribution $F_k(s) = \Pr[s(x_t \mid x_{<t}) \leq s]$. The 90th-percentile surprisal is $Q_k(0.9) = F_k^{-1}(0.9)$. Since deeper generations concentrate probability mass on high-likelihood tokens, $F_k$ stochastically dominates $F_{k+1}$ in the upper tail, yielding $Q_k(0.9) \geq Q_{k+1}(0.9)$.

\paragraph{Surprisal variance (\texttt{var\_surprisal}).}
The variance of token-level surprisal measures the spread of the surprisal distribution within a text. As recursive generation narrows the output distribution toward its modes, the entropy of $M_k$ decreases with $k$. This entropy reduction manifests as a more uniform (lower-variance) surprisal profile: the model becomes increasingly predictable, and tokens cluster near the expected surprisal $H(M_k)$ rather than span a wide range. Formally,
\begin{equation}
    \text{Var}_k[s(x_t \mid x_{<t})] = \mathbb{E}_k[s^2] - (\mathbb{E}_k[s])^2,
    \label{eq:var_surprisal}
\end{equation}
where expectations are over tokens in texts from $\mathcal{D}_k$. Both the mean surprisal and the second moment decrease as the distribution concentrates, but the second moment decreases faster because it weights extreme values quadratically. The net effect is a decrease in variance.

\begin{proposition}[Monotonicity of Feature Signals]
\label{prop:monotone}
Under the recursive generation framework of Equations~\eqref{eq:train}--\eqref{eq:sample}, both \texttt{p90\_ppl} and \texttt{var\_surprisal} are monotonically non-increasing in depth $k$, i.e.,
\begin{equation}
    Q_{k+1}(0.9) \leq Q_k(0.9), \qquad \text{Var}_{k+1}[s] \leq \text{Var}_k[s], \qquad \forall\, k \geq 0.
\end{equation}
\end{proposition}

Our experimental Spearman correlations of $\rho \in [-0.89, -0.82]$ between these features and generation depth across three model families (Pythia-1.4B, GPT-2 XL, OLMo-1B) corroborate this monotonicity prediction.

\subsection{Upper Bound on Effective Detection Depth}
\label{sec:theory:edd}

While Proposition~\ref{prop:monotone} establishes that the feature signal decreases with depth, it does not address whether the signal remains detectable. We now derive an upper bound on the Effective Detection Depth (EDD), defined as the maximum depth $k^*$ at which a binary classifier can distinguish $\mathcal{D}_k$ from $\mathcal{D}_{k+1}$ with AUC above a threshold $\tau$.

The key insight is that the recursive generation mapping $\mathcal{T}: \mathcal{D}_{k-1} \mapsto \mathcal{D}_k$ is a contraction in distributional distance. Following the framework of \citet{alemohammad2024selfconsuming}, who showed that autophagous loops act as contraction mappings on the space of distributions, we formalize this as follows.

\begin{assumption}[Contraction Property]
\label{asm:contraction}
There exists a constant $\gamma \in (0, 1)$ such that for all $k \geq 1$,
\begin{equation}
    d(\mathcal{D}_k, \mathcal{D}_{k+1}) \leq \gamma \cdot d(\mathcal{D}_{k-1}, \mathcal{D}_k),
    \label{eq:contraction}
\end{equation}
where $d(\cdot, \cdot)$ is a distributional distance (e.g., Wasserstein-2 or total variation).
\end{assumption}

This assumption is supported by theoretical results: \citet{shumailov2024collapse} showed that in the Gaussian case, the distributional distance between consecutive iterations decreases geometrically, and \citet{alemohammad2024selfconsuming} demonstrated analogous behavior for generative image models under fully synthetic training loops.

\begin{theorem}[EDD Upper Bound]
\label{thm:edd}
Let $\delta > 0$ be the minimum distributional distance at which a classifier with capacity $\mathcal{C}$ achieves AUC $\geq \tau$. Under Assumption~\ref{asm:contraction}, the Effective Detection Depth satisfies
\begin{equation}
    \text{EDD}(\tau) \leq \left\lfloor \frac{\log(\delta / d_0)}{\log \gamma} \right\rfloor + 1,
    \label{eq:edd_bound}
\end{equation}
where $d_0 = d(\mathcal{D}_0, \mathcal{D}_1)$ is the initial distributional gap.
\end{theorem}

\begin{proof}
By iterating Equation~\eqref{eq:contraction}, the distributional distance at step $k$ satisfies
\begin{equation}
    d(\mathcal{D}_k, \mathcal{D}_{k+1}) \leq \gamma^k \cdot d_0.
\end{equation}
The classifier fails when this distance falls below $\delta$, i.e., when $\gamma^k \cdot d_0 < \delta$. Solving for $k$ gives $k > \log(\delta / d_0) / \log \gamma$. Since $\gamma \in (0,1)$, $\log \gamma < 0$, and noting that $\delta < d_0$ (otherwise depth 0 is already undetectable), the bound follows.
\end{proof}

Theorem~\ref{thm:edd} reveals that EDD depends on three factors: (i) the initial distributional gap $d_0$ between human and first-generation text, (ii) the contraction rate $\gamma$, determined by model capacity and sampling parameters, and (iii) the classifier sensitivity $\delta$, which depends on the feature set and classifier capacity $\mathcal{C}$.

\paragraph{Interpretation.} The bound is logarithmic in $d_0/\delta$, implying that even large initial gaps yield only modest detection depths. This explains the empirically consistent finding of $\text{EDD} = 4$ across our three model families at threshold $\tau = 0.75$: despite different architectures and scales (1B to 1.5B parameters), the contraction rates are similar enough that the detection horizon converges to a narrow range.

\subsection{Alignment with Experimental Observations}
\label{sec:theory:validation}

We verify the theoretical predictions against three experimental findings.

\paragraph{Non-uniform contraction of pairwise discriminability.}
Theorem~\ref{thm:edd} predicts that the distributional distance between adjacent depths contracts with increasing depth. Our binary pairwise AUC values follow this pattern: $\text{AUC}(d_1, d_2) = 0.922$, $\text{AUC}(d_2, d_3) = 0.802$, $\text{AUC}(d_3, d_4) = 0.712$, $\text{AUC}(d_4, d_5) = 0.653$. The AUC-excess contraction ratio decreases from 0.80 to 0.69 across successive pairs, indicating accelerating convergence toward indistinguishability rather than a constant contraction factor.

\paragraph{Consistent EDD across models.}
The EDD bound in Equation~\eqref{eq:edd_bound} depends on $d_0$, $\gamma$, and $\delta$, but not directly on model architecture. If different model families produce similar contraction rates (as predicted by the model collapse theory for models of comparable capacity), they should yield similar EDD values. Our experiments confirm this: Pythia-1.4B, GPT-2 XL, and OLMo-1B all reach $\text{EDD} = 4$ at $\tau = 0.75$.

\paragraph{Feature monotonicity.}
Proposition~\ref{prop:monotone} predicts monotonically decreasing \texttt{p90\_ppl} and \texttt{var\_surprisal}. The observed Spearman correlations of $\rho = -0.82$ to $-0.89$ confirm strong monotonicity, with the deviation from $\rho = -1.0$ attributable to finite-sample noise and text-level variability within each depth.

In summary, the contraction mapping framework provides a unified explanation for the three central empirical phenomena: monotonic feature decay, non-uniform loss of discriminability, and a finite detection horizon that is consistent across model families.
